% ============================================================================
\clearpage
\section*{v4.10.0-r45 Source-Only: Audit of the Mediator Chain and the Fork Behind It}
\addcontentsline{toc}{section}{v4.10.0-r45 Source-Only: Audit of the Mediator Chain}
\label{sec:v4100-r45-mediator-audit}
% ============================================================================

\paragraph{Purpose.}
Layers v4.10.0-r33 to r44 search for a mediator behind the directed term of the OP5 comparison model.  Each layer removes one set structure
and introduces another.  At the author's request this layer draws up the balance, traces the directed term to its origin, and re-evaluates the
chain one to two levels higher.  It is an audit without a decision criterion; its one computation has a result known by hand beforehand and
is not blind.  Earlier layers are not rewritten; Remark r45.2 corrects readings of r39.3, r40.5 and r41.4.

\noindent\textbf{Source-only status.}
This v4.10.0-r45 layer is source-only.  It does not build a PDF, does not assign a new DOI,
and does not authorize a public release.  The v5.3.0 Version DOI
\href{https://doi.org/10.5281/zenodo.24434403}{10.5281/zenodo.24434403} and the Concept DOI
\href{https://doi.org/10.5281/zenodo.19210728}{10.5281/zenodo.19210728} are unchanged by it.

\subsection*{v4.10.0-r45 origin of the directed term}
\label{subsec:v4100-r45-origin}

The numerator \(N_m=\tfrac23(x^2-a_\perp x+a_\perp^2)\) first appears in the notes on the bifocal response (20 November 2026).  There a source is built by
hand from the coupling \(\mathcal L=-hP^{ij}R^{(1)}_{0i0j}\) with \(P=\operatorname{diag}(1,1,0)\), namely \(S^{00}=-a_\perp\), \(S^{0i}=-\omega(Pk)_i\), \(S^{ij}=-xP^{ij}\) per unit
amplitude, and contracted with the Fierz--Pauli numerator; the notes state that this second coupling is set.  The slowed propagation \(r_m\) and the
values \(r_m=\tfrac12\), \(g_m^2=\tfrac3{100}\) are recorded there as comparison inputs, not derived; r32.0 lists the drive operator as set.

\subsection*{v4.10.0-r45 results}
\label{subsec:v4100-r45-results}

\paragraph{Theoremlet v4.10.0-r45.1 (the directed term is an ordinary Fierz--Pauli exchange).}
The source \(S^{\mu\nu}\) above is conserved, \(p_\mu S^{\mu\nu}=0\), and the Lorentz-covariant Fierz--Pauli numerator is
\[
 S_{\mu\nu}S^{\mu\nu}-\tfrac13\,(S^\mu{}_\mu)^2=\tfrac23\,(x^2-a_\perp x+a_\perp^2)=N_m .
\]
The time--transverse components contribute \(-2S_{0i}S^{0i}=-2a_\perp x\); without them the numerator would be \(\tfrac23(a_\perp+x)^2\).
\emph{Proof.} \(S_{\mu\nu}S^{\mu\nu}=a_\perp^2-2a_\perp x+2x^2\) and \(S^\mu{}_\mu=2x-a_\perp\); checked exactly (script \texttt{fp\_\allowbreak directed\_\allowbreak source.wl}).  \(\square\)

\paragraph{Remark v4.10.0-r45.2 (correction of readings of r39--r41).}
In the original construction the mediator is an ordinary massive spin-2 field with covariant numerator, and the direction sits in the coupling:
in the frame \(u\) (the index 0 of \(R_{0i0j}\)) and in the axis (\(P=1-nn\)).  Layers r38 and r39 put the direction into the mediator instead (a tensor
transverse to the axis, a Euclidean norm on the subspace, a coupling to \(\partial_a\partial_b\eta-\eta_{ab}\Box_\parallel\eta\)).  Only that decomposition raises the
question why the time--transverse components do not couple; in the original construction they couple and are needed (r45.1).  The chain
r39.3, r40.5, r41.4 and r42--r44 (decoupling, equal orientation weights, time reversal against reflection) therefore answers a question that arises
in the decomposition of r38/r39 only.  Its computations stand; its weight for OP5 is smaller than read there.

\paragraph{Remark v4.10.0-r45.3 (balance of set structures).}
Over r33--r44 the number of free choices stays at one; each layer moves it: numerator (r35), frame and axis (r36), tension law (r37), scalar (r38),
decoupling (r39), orientation weights (r40), time reversal against reflection (r41), undecided (r42--r44).  What remains set are three inputs of the
comparison model: the directed coupling with frame \(u\) and axis \(n\), the slowed propagation \(r_m\), and the values.  All three are data of a
distinguished frame.

\paragraph{Remark v4.10.0-r45.4 (one level higher: a fork).}
The four recurring candidates, frame and axis of a background (r36), a string between M and T (r37, deferred idea I1), the time orientation of
H118 (r40) and the counterflow of the seed (r41), are four ways of supplying a frame \(u\), an axis \(n\) and a time orientation.  The search circles one
question, where these come from, because one branch of a fork was taken without saying so.  Branch (H), the field side, treats them as properties
of a field or medium; r35--r44 followed it.  Branch (B), the readout side, treats them as data of the readout (observer frame, measuring axis, time
orientation of the readout); the directed term was originally built this way, the observer reading of OP1 is the consistent one of its two
readings, and H118 calls the octant sign a time orientation of the readout.  OP5 itself is the question of the cosmological constant and the
vacuum balance; the comparison model is a local auxiliary construction.  Which branch is pursued is the author's decision.

\subsection*{v4.10.0-r45 frames}
\label{subsec:v4100-r45-frames}

The source and the exchange are in the matter frame M with the frame \(u\) and axis \(n\) of the comparison model.  No quantity of another frame is compared.

\subsection*{v4.10.0-r45 register entries}
\label{subsec:v4100-r45-register}

\begin{center}
\small
\begin{tabular}{@{}>{\raggedright\arraybackslash}p{0.34\textwidth}>{\raggedright\arraybackslash}p{0.58\textwidth}@{}}
\toprule
statement & status after v4.10.0-r45\\
\midrule
directed term \(N_m\) & ordinary Fierz--Pauli exchange of a set directed, conserved source (r45.1)\\
decoupling of \(H_{0i}\) (r39--r44) & arises in the decomposition of r38/r39 only (r45.2)\\
set inputs of the comparison model & directed coupling (\(u\), \(n\)), slowed propagation \(r_m\), values\\
fork & field side (H) against readout side (B); author's decision\\
OP5 & open\\
\bottomrule
\end{tabular}
\end{center}

\subsection*{v4.10.0-r45 notes}
\label{subsec:v4100-r45-notes}

\texttt{notes/\allowbreak research/\allowbreak op5\_\allowbreak mediator\_\allowbreak audit\_\allowbreak 2026\_\allowbreak 11\_\allowbreak 23/} (audit, script \texttt{fp\_\allowbreak directed\_\allowbreak source.wl}, output in
\texttt{run/}).

\subsection*{v4.10.0-r45 ledger and guard update}
\label{subsec:v4100-r45-ledger-guard}

\[
R_{\rm chk}^{v4100\text{-}r45}=(1,1,0)
\]
for (status recorded; computation checked in the cited script; the origin of frame, axis and slowed propagation not derived).

The layer spends no guard:
\[
R_{\rm glob}^{(v4100\text{-}r45,+)}=(0,1,1,1,1),
\qquad
\Gamma_{\rm guard}^{v4100\text{-}r45}=(1,1,1,1,1),
\qquad
\beta_4=1/256\ \text{unspent}.
\]
The layer does not imply exact finite \(C5\) zero closure, exact all-readout terminality, \(C6/T2\), a global photon-clock/RH theorem,
or global all-readout closure.  OP1--OP5, C5 and C6/T2 remain open.
